% Chapter VI -- Summary of Findings, Conclusions, Recommendations (rev L-104)
% Includes the audit-grounded Limitations & Threats to Validity section ---
% the highest-leverage defense asset. Structure mirrors the validity audit
% of 2026-07-25 so every threat names its artifact, quantification,
% mitigation, and residual risk -- nothing hidden for the panel to discover.
% REV L-104: T11 residual aligned with firewall policy (prohibition lowers
% risk; it does not eliminate it); T10 completed to match the stated
% evidence/mitigation/residual pattern; findings gated with VERIFY tokens;
% practice recommendation hedged behind domain-specific validation.

\chapter{Summary of Findings, Conclusions, and Recommendations}

\section{Summary of Findings}
Restating by research question, with pointers rather than restatements of
statistics (full numbers: Chapter~5). [[VERIFY-AGAINST: each finding below
must match the corresponding Chapter 5 table post-transcription; delete
token when confirmed.]]
\begin{enumerate}
\item \textbf{RQ1 (transfer)}: surrogate-trained detection transferred to
both external benchmark strata and consent-gated human text at usable
operating points, with the transfer cost quantified in
Section~5.3.
\item \textbf{RQ2 (paired comparison)}: the principal system exceeded
zero-shot and embedding baselines on the majority of prespecified strata;
exceptions are enumerated with intervals in Section~5.4.
\item \textbf{RQ3 (calibration)}: reliability degraded under distribution
shift in a consistent direction, supporting flagging rather than
automated-threshold use (Section~5.5).
\item \textbf{RQ4 (error localization)}: headroom concentrated in specific
length bands and mixed-authorship slices; evasion sensitivity was bounded
but nonzero; repair attempts closed part of the gap (Section~5.6).
\end{enumerate}

\section{Conclusions}
(1) Training on controlled synthetic proxies is a viable substitute for
scraped real-model corpora that avoids terms-of-service exposure by
construction, \emph{provided} the resulting surrogate--real gap is
measured rather than assumed --- which this study operationalized via its
external tier. (2) Ranking quality and trustworthy confidence dissociate
under shift; evaluation protocols for MGT detection should therefore
report calibration alongside AUROC, a recommendation this thesis advances
to benchmark practice. (3) Slice-resolved auditing (length, authorship,
attack) changes the verdict relative to aggregate reporting; the audit
suite built here is reusable methodological infrastructure. (4) The
study's own governance instrumentation --- frozen gates, provenance
ledgers, dual reporting --- is offered as a template for first-party MGT
research under ethics constraints.

\section{Limitations and Threats to Validity}
This section is grounded clause-for-clause in the project's validity audit
of 2026-07-25 (executive summary: Appendix~G); each threat names its
artifact, quantification, mitigation attempted, and residual risk.
Self-reporting these before the panel asks is a deliberate defense posture.

\subsection{Construct Validity}
\textbf{T1 --- Surrogate--real gap.} Detectors trained on surrogate\_v2 may
key on generator idiosyncrasies absent in wild models. \emph{Evidence}:
tier-asymmetry in Section~5.3; \emph{mitigation}: multi-family generation
grid + external tier; \emph{residual}: gap nonzero by design; magnitude
[[FILL]].

\textbf{T2 --- Construct drift between proposal and execution.} Objectives
were written before results existed; misalignment would undermine the
construct the thesis claims to test. \emph{Evidence}: alignment-matrix
diff of proposal commitments versus deliverables; \emph{mitigation}: the
alignment pass reconciled objectives to deliverables item-by-item and
froze the revised RQ wording; \emph{residual}: low, documented.

\subsection{Internal Validity}
\textbf{T3 --- Mixed authorship confound.} Documents combining human and
machine text behave as neither class and depress both error types.
\emph{Evidence}: mixed\_authorship breakdown; \emph{mitigation}: explicit
stratification and separate reporting, never silent pooling;
\emph{residual}: moderate within contaminated slices.

\textbf{T4 --- Length/headroom confound.} Aggregate metrics mask
band-specific failure. \emph{Evidence}: headroom\_by\_length;
\emph{mitigation}: slice-resolved reporting plus repair attempts;
\emph{residual}: short-band gap persists [[FILL]].

\textbf{T5 --- Salvage-pathway selection.} phd\_rescued shards entered via
a non-prospective route and could select atypical writers. \emph{Evidence}:
chain-of-custody record; \emph{mitigation}: quarantine tier, dual
reporting of all headlines with/without F2; \emph{residual}: bounded ---
no headline flips [[FILL: confirm from dual tables]].

\textbf{T6 --- Analytical flexibility.} Many strata invite cherry-picking.
\emph{Evidence}: stratum count versus prespecification documents;
\emph{mitigation}: frozen feasibility gate, prespecified comparison family
with Holm correction, pilot-frozen analysis code; \emph{residual}: low.

\subsection{External Validity}
\textbf{T7 --- RAID subsampling.} Conclusions generalize to the sampled
strata, not all of RAID; allocation and floors documented in Appendix~D.
\emph{Residual}: moderate; exact reconstitution possible.

\textbf{T8 --- Domain, language, and source coverage.} Paraphrase/domain
breakdowns show uneven coverage, and tangential socio-academic sources
could ground only motivational, not evidentiary, claims. \emph{Evidence}:
stratum coverage table; \emph{mitigation}: scoped claims language
throughout; \emph{residual}: moderate.

\textbf{T9 --- Temporal drift.} Generator families evolve faster than
publication cycles; results characterize the audited snapshot.
\emph{Residual}: unavoidable; stated on every generalization claim.

\subsection{Statistical Conclusion Validity}
\textbf{T10 --- Multiplicity and power.} Many strata inflate false-positive
risk and small cells yield imprecise estimates. \emph{Evidence}: paired
delta interval widths per stratum; \emph{mitigation}: prespecified
comparison families + Holm control + paired bootstrap intervals;
\emph{residual}: wide intervals on small strata; those cells are
interpreted conservatively and never carry headline claims.

\subsection{Integrity Threat (Researcher-Level)}
\textbf{T11 --- Tooling circularity.} The subject of study (detector
evasion) borders on tooling (humanizer pipelines adjacent to the lineage
disclosed in Appendix~H) that must never touch manuscript text; doing so
would create evidentiary circularity and misconduct exposure.
\emph{Mitigation}: binding integrity firewall policy, composition-
provenance ledger shipped as Appendix~H, proactive adviser disclosure;
\emph{residual}: low but nonzero --- prohibition minimizes rather than
eliminates risk; the ledger and disclosure obligations provide ongoing
monitoring.

\section{Recommendations}
\textbf{For practice}: MGT detectors should be evaluated and deployed as
calibrated triage aids, not automated adjudicators; domain-specific
validation should precede any operational reliance, and procurement
evaluations should require slice reports (length, authorship, attack).
\textbf{For research}: (1) benchmark suites should ship calibration and
slice audits natively; (2) surrogate-training recipes need theory for when
proxy-generator diversity closes surrogate--real gaps; (3) first-party
corpus studies should adopt frozen-gate governance as default; (4) repair
of short-band headroom is the highest-value open problem surfaced here.
\textbf{For the institution}: adopt an explicit policy separating research
use of evasion tooling from authoring workflows, per the shipped firewall
policy (Appendix~H).
\textbf{Future work}: extend tiers longitudinally to measure drift;
expand own corpus past gate ceilings to power subpopulation analyses
currently underpowered; formalize intrinsic-dimension probes from
descriptive lens to tested feature family.